\chapter{Manejo en tierra, servicio y mantenimiento por el piloto}
\label{cap:manejo}

\opening{Todo límite, procedimiento, intervalo y requisito de servicio de este capítulo es obligatorio.}

\section{Generalidades}
\fuente{POH, sección 8}

Este capítulo describe los procedimientos recomendados para manejar, dar servicio y mantener el R22. Todo propietario debe mantenerse en estrecho contacto con un Centro de Servicio Robinson para obtener la información más reciente, y registrarse en la fábrica para recibir boletines de servicio, cambios del manual y otra información útil (disponibles en \texttt{www.robinsonheli.com}).

Los reglamentos federales ponen la responsabilidad del mantenimiento en el propietario y el operador, que deben asegurar que todo el mantenimiento sea efectuado por mecánicos calificados y de acuerdo con el Manual de Mantenimiento del R22 (Instrucciones de Aeronavegabilidad Continuada), los boletines y cartas de servicio y las directivas de aeronavegabilidad de la FAA. Todos los límites, procedimientos, prácticas de seguridad, límites de tiempo y requisitos de servicio y mantenimiento del POH se consideran obligatorios.

\section{Documentos requeridos}

El certificado de aeronavegabilidad (formulario FAA 8100-2) debe exhibirse en la aeronave en todo momento. Además deben llevarse a bordo:

\begin{pasos}
\item Certificado de matrícula (formulario FAA 8050-3).
\item Manual de Operación del Piloto.
\item Peso y balance vigente.
\end{pasos}

No deben llevarse a bordo, pero deben estar disponibles para cualquier mecánico o piloto que atienda la aeronave: la bitácora de la aeronave y la bitácora del motor.

\begin{nota}
Los documentos requeridos pueden variar en países distintos de los EE.\,UU.\sidenote{Nota del compilador: en Chile, los documentos exigibles a bordo los fija la normativa DAR de la DGAC.}
\end{nota}

\section{Inspecciones requeridas}

Los reglamentos federales exigen que la mayoría de las aeronaves civiles de matrícula estadounidense se sometan a una inspección completa cada doce meses (inspección anual), firmada por un mecánico con autorización de inspección (IA). Además, el Manual de Mantenimiento del R22 exige una inspección completa cada 100 horas de operación.

El R22 incluye muchas características únicas. Sin entrenamiento especial, los mecánicos de célula y motor (A\&P) no están calificados para efectuar estas inspecciones. Solo el personal debidamente habilitado que haya aprobado un curso de mantenimiento aprobado por la fábrica, o que trabaje bajo su supervisión directa, debe efectuar mantenimiento, reparaciones o inspecciones en el R22.

El helicóptero incorpora varios componentes con vida limitada por fatiga, que deben retirarse a intervalos especificados. Su lista está en la sección de Limitaciones de Aeronavegabilidad del Manual de Mantenimiento (capítulo~\ref{cap:mm-vida-limitada} de este libro).

La fábrica publica ocasionalmente boletines de servicio y la FAA directivas de aeronavegabilidad (AD) aplicables a grupos específicos de aeronaves. Son cambios o inspecciones obligatorias que deben cumplirse dentro del plazo especificado. Los propietarios deben verificar periódicamente con los Centros de Servicio Robinson que se hayan cumplido los últimos boletines y AD.

\section{Mantenimiento preventivo por el piloto}

El 14 CFR parte 43 permite a un piloto certificado efectuar mantenimiento preventivo. Aplicado al R22, incluye:

\begin{pasos}
\item Reemplazar alambre de frenado o pasadores de chaveta defectuosos.
\item Reemplazar ampolletas, reflectores y lentes de las luces de posición y aterrizaje.
\item Reemplazar, limpiar o calibrar la separación de las bujías.
\item Limpiar o reemplazar el filtro de aire del motor.
\item Cambiar el aceite y el filtro del motor.
\item Inspeccionar y limpiar los detectores de virutas.
\item Reponer el aceite de las cajas principal y de cola.
\item Retirar y reemplazar la cuba del gascolator.
\item Retirar o reemplazar capós o paneles de inspección.
\item Dar servicio o reemplazar la batería.
\item Reemplazar las zapatas de desgaste de los patines.
\item Limpiar o repintar el exterior de la aeronave.
\end{pasos}

Aunque la ley lo permite, este trabajo solo debe efectuarlo un piloto seguro de estar calificado para completarlo de manera confiable, y siempre conforme al Manual de Mantenimiento. Al terminar, cuando corresponda, el piloto debe anotar en la bitácora: la fecha, la descripción del trabajo, las horas totales de la aeronave, el número de su licencia y su firma.

\section{Alteraciones de la aeronave}

Lo compacto del R22 y sus muchas características de diseño únicas hacen desaconsejable cualquier modificación. Las características dinámicas y la susceptibilidad a la fatiga de los sistemas rotor, de transmisión y de control hacen que cualquier modificación de ellos sea extremadamente peligrosa.

También es peligroso instalar cualquier equipo electrónico o aviónica no aprobada ni suministrada por la fábrica. Por lo compacto de la consola y del túnel que contiene los controles y los mazos de cables, es probable que cualquier cable adicional interfiera con el libre movimiento de los controles. Los tacómetros electrónicos y otros instrumentos pueden verse afectados por otros dispositivos eléctricos, y su confiabilidad y exactitud son esenciales para una operación segura. Por estos riesgos, Robinson no autoriza ninguna modificación o alteración que no sea suministrada por la fábrica e instalada por personal entrenado por ella.

\section{Manejo en tierra}

Para nivelar, izar o levantar con gatas, vea el Manual de Mantenimiento (capítulo~\ref{cap:mm-datos-servicio}).

El helicóptero puede moverse en tierra con las ruedas de maniobra. La presión recomendada de los neumáticos es 60~psi (4~bar). Las ruedas deben retirarse para volar. Para instalarlas:

\begin{pasos}
\item Sostenga la manilla y la rueda con el eje saliente en su posición más baja.
\item Inserte el eje en el soporte montado en el patín, asegurándose de que entre por completo.
\item Lleve la manilla más allá del punto muerto para levantar el helicóptero y trabar la rueda en posición.
\end{pasos}

\begin{precaucion}
Al bajar el helicóptero, la manilla tiende a saltar bruscamente.
\end{precaucion}

Mueva el helicóptero sujetando la caja de cola y la sección trasera del cono de cola. Si se necesita ayuda, una segunda persona puede empujar uno de los tubos verticales traseros del bastidor o la nariz. También puede usarse un carro de remolque eléctrico Robinson según sus instrucciones.

\begin{precaucion}
No mueva el helicóptero tomándolo de los estabilizadores verticales, la parte exterior del estabilizador horizontal, el rotor de cola, el patín de cola ni los mandos del rotor de cola.
\end{precaucion}

\section{Estacionamiento}

\begin{pasos}
\item Coloque el cíclico en neutro y aplique fricción.
\item Ponga el colectivo completamente abajo y aplique fricción.
\item Alinee las palas aproximadamente de proa a popa. Aplique el freno rotor. Use amarras de palas con viento.
\end{pasos}

\begin{precaucion}
Si usa amarras de palas, no apriete en exceso las correas (tensión máxima 5~lb). No tire de las palas hacia abajo para balancear el rotor: para bajar una pala, empuje hacia arriba la opuesta.
\end{precaucion}

\begin{pasos}
\setcounter{enumi}{3}
\item Con tormenta, el helicóptero debe guardarse en hangar o trasladarse a un área segura.
\item Si se guarda en hangar, cierre la válvula de corte de combustible.
\end{pasos}

\section{Puertas de la cabina}

Ambas puertas pueden ser retiradas e instaladas por el personal de mantenimiento o los pilotos. Para retirar una puerta, saque los pasadores de chaveta de los pasadores de bisagra superior e inferior, ábrala y levántela. Para instalarla, invierta el procedimiento. Ajuste el peso y balance según corresponda al retirar o instalar las puertas.

\section{Aceite del motor}

La cantidad máxima recomendada es de seis cuartos y la mínima para despegue de cuatro cuartos, según la varilla.

\textbf{El aceite y el filtro deben cambiarse al menos cada 50 horas o cuatro meses, lo que ocurra primero. Si no hay filtro de aceite instalado, cambie el aceite cada 25 horas.} Para cambiar el aceite:

\begin{pasos}
\item Haga funcionar el motor en tierra o vuele para alcanzar la temperatura normal de operación.
\item Retire los faldones laterales izquierdos. Si hay soplador de calefacción, desconecte la manguera de calefacción y el cable de masa del motor del soplador y retire el perno que fija el motor del soplador al bastidor. Abra el drenaje rápido bajo el cárter para drenar el aceite en un recipiente adecuado.
\item Corte el alambre de frenado de la tapa del filtro de succión de aceite, cerca del eje central en el extremo de accesorios (magnetos) del cárter (en los motores O-320, está en la parte inferior del cárter).
\item Retire, inspeccione, limpie y reinstale el filtro de succión. Vuelva a frenar la tapa.
\item Corte el alambre de frenado del filtro de aceite (si lo tiene) y aflójelo con una llave, o afloje los pernos de la carcasa del filtro de presión.
\item Coloque un recipiente adecuado debajo y hacia el interior del magneto para recoger el aceite retenido en el filtro o en la carcasa del filtro, y retire el filtro o la carcasa lentamente. No deje que el aceite escurra sobre la carcasa del magneto.
\item Corte el filtro para inspeccionarlo, o inspeccione y limpie el filtro de presión.
\item Instale un filtro nuevo según las instrucciones impresas en él. Use solo filtros Champion CH48108, CH48108-1, Tempest AA48108-2 o Robinson B123-1 o -2. El filtro no requiere alambre de frenado. Reinstale el filtro de presión y su carcasa, y apriete los pernos de la carcasa a 96~in-lb.
\item Cierre el drenaje rápido y llene el cárter con 6 cuartos del grado de aceite apropiado.
\item Arranque el motor. Verifique presión de aceite dentro de 30 segundos. Hágalo funcionar en tierra unos minutos, deténgalo y verifique que no haya fugas.
\item Revise el nivel con la varilla.
\item Reinstale los faldones laterales y, si corresponde, el motor del soplador y la manguera de calefacción.
\item Haga las anotaciones de mantenimiento correspondientes.
\end{pasos}

\begin{table}[htbp]
\centering\small
\caption{Grados de aceite de motor recomendados.}
\label{tab:aceite-motor}
\begin{tabularx}{\linewidth}{@{}lYY@{}}
\toprule
Temperatura ambiente media & Mineral (MIL-L-6082 o SAE J1966) --- primeras 50 h & Dispersante sin cenizas (MIL-L-22851 o SAE J1899) --- después de 50 h \\
\midrule
Todas las temperaturas & --- & SAE 15W50 o 20W50 \\
Sobre 60~°F (16~°C) & SAE 50 & SAE 40 o 50 \\
30 a 90~°F ($-1$ a 32~°C) & SAE 40 & SAE 40 \\
0 a 70~°F ($-18$ a 21~°C) & SAE 30 & SAE 30, 40 o 20W40 \\
0 a 90~°F ($-18$ a 32~°C) & SAE 20W50 & SAE 20W50 o 15W50 \\
Bajo 10~°F ($-12$~°C) & SAE 20 & SAE 30 o 20W30 \\
\bottomrule
\end{tabularx}
\end{table}

\section{Aceite de las cajas de transmisión}

La cantidad de aceite de las cajas principal y de cola se indica en visores. El visor de la caja principal está en su costado derecho y se ve abriendo la puerta del capó bajo el tanque auxiliar. El de la caja de cola está en su extremo trasero.

Si con el helicóptero nivelado no se ve aceite en el visor, debe agregarse. Las cajas usan aceite Robinson P/N A257-2 o A257-22; consulte el letrero para asegurarse de agregar el aceite correcto. Para agregar aceite a la caja de cola:

\begin{pasos}
\item Corte el alambre de frenado (si lo tiene) y retire la tapa de llenado/respiradero en la parte superior de la caja.
\item Llene muy lentamente hasta el centro del visor. \textbf{No llene en exceso} (normalmente se requiere menos de una cucharadita de aceite).
\item Reinstale la tapa, verificando que la junta metálica o el anillo O esté en su lugar. Apriete a 100~in-lb.
\item Frene con alambre como estaba (no se requiere si la tapa sella con anillo O), asegurándose de que el alambre tire en el sentido que apretaría la tapa.
\end{pasos}

Para agregar aceite a la caja principal, vea el Manual de Mantenimiento.

\section{Combustible}

Los grados y la capacidad de combustible aprobados están en el capítulo~\ref{cap:limitaciones}.

Antes del primer vuelo del día y después de cargar combustible, drene una pequeña cantidad del gascolator y de cada tanque con los drenajes rápidos. En los helicópteros posteriores, un tubo de vidrio guardado dentro de la puerta superior trasera del capó sirve para recoger las muestras. Drene lo suficiente para eliminar agua o suciedad y verifique el color del combustible aprobado. Si se sospecha contaminación, siga drenando del gascolator y de los tanques hasta eliminarla por completo.

\section{Vejigas de combustible}

Los tanques pueden tener vejigas flexibles resistentes al impacto. Las vejigas se deterioran con el tiempo, lo que puede provocar fugas o rezumado de combustible. Para prolongar su vida, llene los tanques al guardar el helicóptero, de modo que el material no se seque. El calor extremo y el agua dañan el material: asegúrese de drenar el agua de los sumideros y guarde el helicóptero en hangar cuando sea posible.

\section{Extintor de incendios (opcional)}

El extintor opcional (modelo RT-A400) debe pesarse mensualmente para verificar que contiene carga suficiente de agente. Anote el resultado en la etiqueta del extintor o en otros registros de mantenimiento. Si el peso baja de 500~gramos, retire el extintor de servicio.

El extintor está clasificado para una temperatura de almacenamiento y operación de 120~°F (49~°C). Si las condiciones lo expondrían a temperaturas mayores (por ejemplo, estacionado al sol en clima caluroso), retírelo y guárdelo en un lugar más fresco entre vuelos.

\section{Batería}

La batería de 12~V está en el compartimiento del motor o bajo la consola de instrumentos. Es sellada y no requiere revisar el nivel de líquido. \textbf{Una batería descargada NO es aeronavegable}, porque no tendrá la reserva necesaria para alimentar el sistema eléctrico si el sistema de carga falla en vuelo.

A menudo una carga de 10 o 15 minutos mejora lo suficiente la batería para arrancar el motor. Si la batería está en el compartimiento del motor, conecte primero el cable positivo del cargador al terminal positivo (lado batería) del relé de batería y luego el negativo a la cinta de masa de la batería o al motor. Si está bajo la consola, retire los tornillos de sujeción de la consola (uno por lado), levántela, retire la tapa de la caja de batería y conecte los cables directamente a los bornes (primero el positivo). Tenga extremo cuidado de no hacer cortocircuito con la chapa de la consola.

Después de cargar, desconecte los cables (primero el negativo), asegure la consola si la abrió e intente un arranque normal. Si la batería sigue sin carga suficiente para arrancar, dele servicio o reemplácela antes del siguiente vuelo.

\section{Limpieza del helicóptero}

\subsection{Exterior, incluidas las palas}

Lave el helicóptero con jabón suave y agua. No use abrasivos fuertes, jabones alcalinos ni detergentes: pueden rayar las superficies pintadas o plásticas o corroer el metal. Cubra las áreas que la solución de limpieza pueda dañar.

\begin{pasos}
\item Enjuague con agua la suciedad suelta y los residuos de sal.
\item Aplique la solución de limpieza con un paño suave, una esponja o un cepillo de cerdas suaves.
\item Para el aceite y la grasa persistentes, use un paño humedecido con nafta alifática.
\item Enjuague a fondo todas las superficies.
\item Aplique cera de carnaúba a las palas y renuévela cuando el agua deje de formar gotas en su superficie. Para las demás superficies pintadas sirve cualquier buena cera automotriz. Use paños suaves o gamuza para no rayar al limpiar o pulir.
\end{pasos}

\begin{precaucion}
Nunca use agua a alta presión para limpiar el helicóptero. Nunca sople aire comprimido en los agujeros de drenaje de las puntas de las palas del rotor principal o de cola.
\end{precaucion}

\subsection{Parabrisas y ventanas}

\begin{pasos}
\item Retire con agua limpia la suciedad, el barro y otras partículas sueltas de las superficies exteriores.
\item Lave con jabón suave y agua tibia o con limpiador de plásticos de aviación. Use un paño suave o esponja con movimiento recto de ida y vuelta, sin frotar con fuerza.
\item Elimine el aceite y la grasa con un paño humedecido en alcohol isopropílico o nafta alifática.
\end{pasos}

\begin{precaucion}
No use gasolina, otros alcoholes, benceno, tetracloruro de carbono, diluyente, acetona ni limpiavidrios en aerosol.
\end{precaucion}

\begin{pasos}
\setcounter{enumi}{3}
\item Después de limpiar las superficies plásticas, aplique una capa delgada de cera de pulir dura y frote suavemente con un paño suave, sin movimientos circulares.
\end{pasos}

\begin{precaucion}
La superficie del parabrisas debe ser hidrófoba (repeler el agua) para tener buena visibilidad con lluvia. Al usar un producto de limpieza o pulido nuevo en el parabrisas, verifique antes de volar que el agua forme gotas en la superficie.
\end{precaucion}

\begin{pasos}
\setcounter{enumi}{4}
\item En las ventanas acrílicas (parabrisas estándar), las rayas pueden eliminarse frotando con rojo de joyero y puliendo luego a mano con pulidor comercial de plásticos, con un movimiento en forma de ocho.
\end{pasos}

\begin{nota}
Los parabrisas resistentes al impacto son de policarbonato con un recubrimiento duro protector y no pueden pulirse.
\end{nota}

\subsection{Tapicería y asientos}

\begin{pasos}
\item Aspire y cepille, y luego pase un paño húmedo. Seque de inmediato.
\item La tapicería sucia, excepto el cuero, puede limpiarse con un buen limpiador apropiado para el material, siguiendo las instrucciones del fabricante. Evite empaparla o frotarla con fuerza.
\end{pasos}

\begin{precaucion}
Si hay detector de CO instalado, evite solventes, detergentes o aerosoles cerca del sensor y cubra con cinta las aberturas superior e inferior de su carcasa al limpiar el interior.
\end{precaucion}

\begin{pasos}
\setcounter{enumi}{2}
\item El cuero debe limpiarse con jabón para monturas o con un jabón duro suave y agua.
\end{pasos}

\subsection{Alfombras}

Retire la suciedad suelta con una escobilla o aspiradora. Para manchas, use un líquido de limpieza en seco no inflamable.

\section{Almacenamiento}

El helicóptero requiere una preparación especial para el almacenamiento prolongado (más de 30 días). Consulte a su proveedor de mantenimiento los procedimientos apropiados antes de almacenarlo (vea también el capítulo~\ref{cap:mm-datos-servicio}).
